\section{Star-connected boundaries and connected thickening}
\label{s:topology}

Two sites are nearest neighbours when their ambient \(\ell^1\)
distance is one, and star-neighbours when their ambient
\(\ell^\infty\) distance is one.  A path in either graph is
a finite sequence of vertices with consecutive vertices adjacent
in that graph.  A one-point set is connected in both senses.

\begin{lemma}\label{l:boundaryconn}
If \(A\subset\mathbb Z^2\) is finite, nonempty and star-connected,
and \(A^c\) is star-connected, then
\(\partial_\infty A=\{z\notin A:\operatorname{dist}_\infty(z,A)=1\}\)
is nonempty and star-connected.
Here star-neighbours differ by \(\ell^\infty\) distance one.
\end{lemma}

\begin{proof}
Let \(\mathcal K\) be the graph whose vertices are \(\mathbb Z^2\)
and whose edges join star-neighbours.  Finite edge sets form
a vector space over \(\mathbb F_2\), with addition given by
symmetric difference.  An edge set is called even if each
vertex has even degree in it.  A triangle consists of three
pairwise star-neighbouring vertices and its boundary is its
three-edge set.  Such a triangle lies in a square
\((i,j)+\{0,1\}^2\).

We first prove that every finite even edge set is a finite
sum of triangle boundaries.  Let \(Z\) be such an edge set.
For each diagonal edge \(e=\{p,p+(1,\epsilon)\}\) of \(Z\),
where \(\epsilon\in\{-1,1\}\), add the boundary of the triangle
\(\{p,p+(1,0),p+(1,\epsilon)\}\).  Its only diagonal edge
is \(e\).  The resulting edge set \(Z'\) is finite, even,
and has horizontal and vertical edges only.

Write \(H(i,j)\in\mathbb F_2\) for the indicator that
\(\{(i,j),(i+1,j)\}\) belongs to \(Z'\), and
\(V(i,j)\) for the analogous indicator of
\(\{(i,j),(i,j+1)\}\).  Choose an integer \(R\ge1\) such that
every endpoint of \(Z'\) lies in \([-R,R]^2\); when \(Z'\) is
empty take \(R=1\).  In particular \(H(i,r)=0\) for
\(|r|>R\), and \(V(i,R)=0\).  Define the finite sum
\[
 n(i,j)=\sum_{\substack{-R\le r\le R\\r>j}}H(i,r)
                    \quad\hbox{in }\mathbb F_2.
\]
For each fixed \(i\), \(\sum_{r=-R}^R H(i,r)=0\).  Indeed this sum
is the parity of the number of edges crossing the line from
the half-plane \(x_1\le i\) to its complement.  It equals
the sum of the degrees of \(Z'\) at vertices of that
half-plane, because edges inside it contribute twice.
This sum is zero because \(Z'\) is even.  All degree sums
are finite after omitting the zero terms, since \(Z'\) is
finite.  It follows that \(n(i,j)\) is nonzero for only
finitely many pairs \((i,j)\): only finitely many columns
contain horizontal edges, and in each such column the sum
vanishes both above its highest horizontal edge and below
its lowest one.

Let \(P_{ij}=(i,j)+\{0,1\}^2\), and let \(\partial P_{ij}\)
be its four nearest-neighbour edges.  We claim that
\[
                    Z'=\sum_{i,j}n(i,j)\partial P_{ij}.
\]
The horizontal edge with indicator \(H(i,j)\) occurs on
the right with coefficient
\(n(i,j-1)+n(i,j)=H(i,j)\).
The vertical edge with indicator \(V(i,j)\) occurs with
coefficient
\[
 n(i-1,j)+n(i,j)
       =\sum_{\substack{-R\le r\le R\\r>j}}
                          \bigl(H(i-1,r)+H(i,r)\bigr).
\]
For \(j<R\), sum the even-degree identity
\(H(i-1,r)+H(i,r)+V(i,r-1)+V(i,r)=0\)
over the finite integer interval \(j+1\le r\le R\).
Terms with \(r<-R\) are zero by the endpoint bound.
All vertical terms cancel in pairs except \(V(i,j)\) and
\(V(i,R)=0\), proving that this coefficient is \(V(i,j)\).
For \(j\ge R\), both the sum and \(V(i,j)\) are zero.
The claim follows.  Finally, each \(\partial P_{ij}\) is
the sum of the two triangle boundaries obtained by inserting
one of the diagonals of \(P_{ij}\).  Together with the
initial diagonal eliminations, this proves the assertion
about even edge sets.

Let \(\mathcal B\) be the set of edges of \(\mathcal K\)
with one endpoint in \(A\) and one in \(A^c\).  It is finite,
and it is nonempty because a point of \(A\) of largest first
coordinate has its right nearest neighbour outside \(A\).
Say that two edges of \(\mathcal B\) are linked if some
triangle boundary contains both.  We prove that any two
edges of \(\mathcal B\) can be joined by a finite sequence
of linked edges.  Otherwise choose a connected component
\(\mathcal B_1\) for this relation and put
\(\mathcal B_2=\mathcal B\setminus\mathcal B_1\).
Both parts are nonempty, and no linked pair crosses between
them.  For finite edge sets define the linear functional
\[
              \varphi(Z)=|Z\cap\mathcal B_1|\pmod2.
\]
A triangle has either zero or two edges in \(\mathcal B\),
according to whether all its vertices lie on the same side
of \(A\), or they do not.  In the latter case the two edges
are linked and belong to the same \(\mathcal B_r\).
Thus \(\varphi\) vanishes on every triangle boundary, and
the preceding decomposition implies that it vanishes on
every finite even edge set.

Choose \(e_r=\{a_r,b_r\}\in\mathcal B_r\), with
\(a_r\in A\) and \(b_r\in A^c\), for \(r=1,2\).
There is a finite star-path from \(a_1\) to \(a_2\) in
\(A\), and a finite star-path from \(b_2\) to \(b_1\)
in \(A^c\), by the two connectivity assumptions.
Concatenate these paths with \(e_2\) and \(e_1\) to form
a closed walk.  The set \(Z\) of edges traversed an odd
number of times by this walk is finite and even: at every
vertex the number of entrances equals the number of exits
modulo two.  Its intersection with \(\mathcal B\) consists
of exactly the distinct edges \(e_1,e_2\), since both
connecting paths stay on one side.  This gives
\(\varphi(Z)=1\), a contradiction.  The linked-edge relation
is therefore connected.

Every edge of \(\mathcal B\) has a unique endpoint in
\(A^c\), and these outer endpoints are precisely the
vertices of \(\partial_\infty A\).  Outer endpoints of
two linked edges lie in the same triangle, so they are
equal or star-neighbours.  A sequence of linked edges
thus gives a star-path between their outer endpoints,
after deleting consecutive repetitions.  Surjectivity of
the outer-endpoint map proves star-connectivity of
\(\partial_\infty A\); nonemptiness follows from that of
\(\mathcal B\).
\end{proof}

\begin{lemma}\label{l:thicken}
If \(X_0\) is finite, nonempty and star-connected, then
\(X=\{z:d(z,X_0)\le1\}\) is nearest-neighbour connected.
For any finite \(G\) containing \(X\), every \(0<\vartheta<1\)
satisfies
\[
 \sum_{z\in X}\vartheta^{d(z,G^c)}
       \le5\vartheta^{-1}
                      \sum_{x\in X_0}\vartheta^{d(x,G^c)}.
\]
\end{lemma}

\begin{proof}
For each \(x\) put
\(C(x)=\{x,x+(1,0),x-(1,0),x+(0,1),x-(0,1)\}\).
This five-point cross is nearest-neighbour connected,
and \(X=\bigcup_{x\in X_0}C(x)\).
For star-neighbours \(x,x'\), if their difference has
one zero coordinate then they are nearest neighbours
and \(C(x)\cap C(x')\) contains both centres.
Otherwise \(x'-x=(\epsilon_1,\epsilon_2)\), where
\(\epsilon_1,\epsilon_2\in\{-1,1\}\), and their crosses
share the two points \(x+(\epsilon_1,0)\) and
\(x+(0,\epsilon_2)\).  These are all possibilities for
star-neighbours.  A star-path between centres therefore
gives a chain of intersecting connected crosses.  Any
two points of their union can be joined by concatenating
paths inside the crosses through these intersection
points.  This proves nearest-neighbour connectedness
of \(X\), including when \(X_0\) is a singleton.

For finite \(G\supset X\), its complement is nonempty
and the ambient distance \(b(z)=d(z,G^c)\) is
one-Lipschitz.  Thus for \(z\in C(x)\),
\(b(z)\ge b(x)-1\).  Since \(0<\vartheta<1\),
\(\vartheta^{b(z)}\le\vartheta^{-1}\vartheta^{b(x)}\).
Every point of \(X\) is in at least one of the crosses,
so summing nonnegative terms gives
\[
 \sum_{z\in X}\vartheta^{b(z)}
    \le\sum_{x\in X_0}\sum_{z\in C(x)}\vartheta^{b(z)}
    \le5\vartheta^{-1}\sum_{x\in X_0}\vartheta^{b(x)}.
\]
This proves the claimed estimate without requiring
the crosses to be disjoint.
\end{proof}
